\clearpage
\begin{summarybox}
\textbf{一句话记住：同样是焦点夹角，椭圆看法线，双曲线看切线；
抛物线先看轴向射线是否指向开口。}
\[
\begin{array}{c|c|c}
\text{曲线}&\text{内部角的两条边}&\text{平分内角的线}\\\hline
\text{椭圆}&PF_1,PF_2&\text{法线}\\
\text{双曲线}&PF_1,PF_2&\text{切线}\\
\text{抛物线}&PF,PH\ (PH\text{ 指向开口})&\text{法线}
\end{array}
\]
\textbf{方法选用：}
最低门槛用“\textbf{对称＋距离和/差}”；
速度最快用“\textbf{角平分线定理}”；
步骤最少用“\textbf{两个单位方向相加}”；
常规证明用“\textbf{求切线＋比较焦点到直线的距离}”。

\textbf{考场检查顺序：}
先定 $P$ 与焦点的射线方向；
再看哪一条线落在夹角内；
最后判断内角还是外角。
双曲线左支要重新核对焦半径大小，
抛物线的 $PH$ 要画箭头；顶点处的角需单独判断。
\end{summarybox}
